% T10 — «Контрольная»: 2 варианта параллельно, графитовый, серьёзный формат с грифом.
% Профиль: математика, квадратные уравнения. 5 задач × 2 варианта.
\documentclass[a4paper,11pt]{article}

\usepackage{../../../../Lessons/_templates/preamble_worksheet}
\usepackage{../_shared/worksheet_check}

% --- Палитра: графитовый, строгий ---
\definecolor{wcprimary}{HTML}{2D3748}
\definecolor{wcprimaryLight}{HTML}{A0AEC0}
\definecolor{wcprimaryBG}{HTML}{EDF2F7}

% --- Заголовок: гриф «КОНТРОЛЬНАЯ РАБОТА» ---
\renewcommand{\WorksheetTitle}[1]{%
  \par\noindent\begin{tikzpicture}
    \draw[wcprimary, line width=1.2pt] (0,0) -- (\linewidth,0);
    \node[anchor=north, text=wcprimary, font=\footnotesize\bfseries\sffamily]
      at (\linewidth/2, -1mm) {\textsc{Контрольная работа}};
    \node[anchor=north, text=wcprimary, font=\fontsize{18pt}{22pt}\selectfont\bfseries]
      at (\linewidth/2, -7mm) {#1};
    \draw[wcprimary, line width=1.2pt] (0,-18mm) -- (\linewidth,-18mm);
  \end{tikzpicture}\par\vspace{3mm}}

% --- TaskBoxVar для двух вариантов параллельно ---
\renewcommand{\TaskBox}[2]{%
  \par\nopagebreak\noindent\begin{tikzpicture}
    \node[anchor=north west, inner xsep=0pt, inner ysep=2pt,
          text width=\linewidth] (B) at (0,0) {%
      {\bfseries\color{wcprimary}№#1.}\ \small #2};
  \end{tikzpicture}\par\vspace{3mm}}

\SetAnswerFieldSize{24mm}{8mm}

\begin{document}

\WorksheetTitle{Квадратные уравнения}
\par\noindent\begin{tabular}{@{}p{0.48\linewidth}@{\hspace{4mm}}p{0.48\linewidth}@{}}
{\small\itshape\color{wcprimary!80}ФИО:\ \rule{30mm}{0.4pt}} & {\small\itshape\color{wcprimary!80}Класс:\ \rule{16mm}{0.4pt}\quad Дата:\ \rule{22mm}{0.4pt}}
\end{tabular}\par\vspace{4mm}

\noindent
\begin{tabular}{@{}p{0.47\linewidth}@{\hspace{6mm}}p{0.47\linewidth}@{}}
% Шапки вариантов
\multicolumn{1}{@{}c@{}}{%
  \begin{tikzpicture}\node[fill=wcprimary, text=white, font=\bfseries\sffamily\small,
    inner xsep=8mm, inner ysep=1.5mm, rounded corners=2pt]{Вариант 1};\end{tikzpicture}} &
\multicolumn{1}{@{}c@{}}{%
  \begin{tikzpicture}\node[fill=wcprimary, text=white, font=\bfseries\sffamily\small,
    inner xsep=8mm, inner ysep=1.5mm, rounded corners=2pt]{Вариант 2};\end{tikzpicture}}\\[3mm]
% Задача 1
{\bfseries\color{wcprimary}№1.}\ Решите уравнение $x^2 - 5x + 6 = 0$. В ответе укажите больший корень.\par\vspace{1mm}\hfill\textbf{Ответ:}~\AnswerField{1} &
{\bfseries\color{wcprimary}№1.}\ Решите уравнение $x^2 - 7x + 12 = 0$. В ответе укажите больший корень.\par\vspace{1mm}\hfill\textbf{Ответ:}~\AnswerField{2}\\[3mm]
% Задача 2
{\bfseries\color{wcprimary}№2.}\ Решите уравнение $2x^2 - 5x - 3 = 0$. Если корней несколько, в ответ запишите их сумму.\par\vspace{1mm}\hfill\textbf{Ответ:}~\AnswerField{3} &
{\bfseries\color{wcprimary}№2.}\ Решите уравнение $3x^2 - 7x - 6 = 0$. Если корней несколько, в ответ запишите их сумму.\par\vspace{1mm}\hfill\textbf{Ответ:}~\AnswerField{4}\\[3mm]
% Задача 3
{\bfseries\color{wcprimary}№3.}\ При каком значении $p$ уравнение $x^2 + px + 9 = 0$ имеет единственный корень? В ответ запишите больший $p$.\par\vspace{1mm}\hfill\textbf{Ответ:}~\AnswerField{5} &
{\bfseries\color{wcprimary}№3.}\ При каком значении $p$ уравнение $x^2 + px + 16 = 0$ имеет единственный корень? В ответ запишите больший $p$.\par\vspace{1mm}\hfill\textbf{Ответ:}~\AnswerField{6}\\[3mm]
% Задача 4
{\bfseries\color{wcprimary}№4.}\ Сумма корней уравнения $x^2 - 8x + k = 0$ равна $8$, а произведение равно $15$. Найдите $k$.\par\vspace{1mm}\hfill\textbf{Ответ:}~\AnswerField{7} &
{\bfseries\color{wcprimary}№4.}\ Сумма корней уравнения $x^2 - 10x + k = 0$ равна $10$, а произведение равно $21$. Найдите $k$.\par\vspace{1mm}\hfill\textbf{Ответ:}~\AnswerField{8}\\[3mm]
% Задача 5
{\bfseries\color{wcprimary}№5.}\ Решите уравнение $x^4 - 5x^2 + 4 = 0$. В ответе укажите наибольший корень.\par\vspace{1mm}\hfill\textbf{Ответ:}~\AnswerField{9} &
{\bfseries\color{wcprimary}№5.}\ Решите уравнение $x^4 - 10x^2 + 9 = 0$. В ответе укажите наибольший корень.\par\vspace{1mm}\hfill\textbf{Ответ:}~\AnswerField{10}\\
\end{tabular}

\vspace{6mm}
{\small\itshape\color{wcprimary!80}Время выполнения: 40 минут. Калькулятор запрещён. Записывайте только окончательные ответы.}

\WorksheetQR{T10-quadratic/L1/v1}

\end{document}
